% ACCEPT -- os/samir/cmd/proc.c p_accept (prompt expression, then one
% cooked line into a Character memvar; empty line -> "").  Gate:
% test-interp-proc.
\begin{XbaseEntry}{ACCEPT}{command}
Displays a prompt, reads a line typed at the keyboard and stores it in a
memory variable as a character string.
\Syntax
\SyntaxLine{\kw{ACCEPT} \opt{\xp{expC}} \kw{TO} \xp{memvar}}
\Usage
\cmd{ACCEPT} displays the character expression \xp{expC}, if one is given, and
waits for a line to be typed.  The line ends when you press \key{Enter}.  The
characters typed are stored in \xp{memvar} as a character string, whatever
they are: a number typed in answer to \cmd{ACCEPT} is stored as characters.  If
you press \key{Enter} alone, \xp{memvar} is an empty string.  The memory
variable is created if it does not exist.

The prompt is usually a character string in quotation marks.  It is displayed
as it stands, so end it with a blank or a colon and a blank to keep the answer
apart from it.
\Tips
Use \cmd{INPUT} when the answer is a number, a date or a logical value, and
\cmd{WAIT} when one keystroke is enough.  To compare an answer with stored
data in capitals, convert it with \fn{UPPER()}.
\Examples
\DotExample{examples/c-accept}
Without a prompt, \cmd{ACCEPT} waits on a new line:
\DotExample{examples/c-accept-2}
\SeeAlso{INPUT, STORE, UPPER(), WAIT}
\end{XbaseEntry}

